\lyseis{Α1.3}{Πρώτοι και σύνθετοι αριθμοί}
\sol{1} α) Σωστό ($24=8\cdot 3$). β) Λάθος (ο $12$ είναι μεγαλύτερος από τον $6$). γ) Σωστό. δ) Λάθος (ο $2$ είναι άρτιος και πρώτος). ε) Σωστό ($121=11\cdot 11$). στ) Λάθος (το $1$ δεν είναι ούτε πρώτος ούτε σύνθετος).
\sol{2} $36$: $1,2,3,4,6,9,12,18,36$. $50$: $1,2,5,10,25,50$ (πρώτοι $2,5$). $83$: μόνο $1$ και $83$, άρα πρώτος.
\sol{3} $2,3,5,7,11,13,17,19,23,29$.
\sol{4} α) πρώτος. β) σύνθετος ($3\cdot 13$). γ) σύνθετος ($3\cdot 17$). δ) σύνθετος ($3\cdot 19$). ε) πρώτος. στ) σύνθετος ($7\cdot 13$).
\sol{5} Ο $120$: $2,3,4,5,10$ (όχι $9$, άθροισμα $3$). Ο $345$: $3,5$ (άθροισμα $12$). Ο $1.008$: $2,3,4,9$ (άθροισμα $9$, $08=8$ διαιρείται με το $4$).
\sol{6} α) $640$, το ψηφίο είναι $0$. β) Το άθροισμα είναι $5+3+8+x=16+x$. Διαιρείται με το $3$ για $x=2,5,8$. Με το $9$ διαιρείται για $x=2$ ($18$). Άρα $x=5$ ή $x=8$. γ) Άθροισμα $1+0+3+4+x=8+x$, άρτιος: $x=4$ ($12$ διαιρείται με το $3$). Άρα $10.344$. δ) Για το $10$ πρέπει $x=0$: $9.700$, και $00$ άρα διαιρείται και με το $4$.
\sol{7} α) Ναι. β) Όχι ($4$ δεν είναι πρώτος). γ) Ναι. δ) Όχι ($9$ δεν είναι πρώτος). ε) Ναι.
\sol{8} $420=2^2\cdot 3\cdot 5\cdot 7$. $792=2^3\cdot 3^2\cdot 11$. $1.350=2\cdot 3^3\cdot 5^2$.
\sol{9} Οι τρόποι αντιστοιχούν στους διαιρέτες: ο $24$ έχει $8$ διαιρέτες ($1,2,3,4,6,8,12,24$), άρα $8$ τρόποι. Ο $13$ είναι πρώτος και έχει $2$ διαιρέτες: μία σειρά των $13$ ή $13$ σειρές του ενός. Δεν υπάρχει άλλη ορθογώνια διάταξη.
\sol{10} Ο $51$ έχει άθροισμα ψηφίων $5+1=6$, που διαιρείται με το $3$, άρα ο $51$ διαιρείται με το $3$: $51=3\cdot 17$. Άρα είναι σύνθετος.
\sol{11} Ο αριθμός είναι $2\cdot2\cdot 3\cdot 5\cdot 7$. Το $6=2\cdot 3$ και το $35=5\cdot 7$ είναι γινόμενα παραγόντων του, άρα τον διαιρούν. Ομοίως $60=2^2\cdot 3\cdot 5$ και $21=3\cdot 7$. Με το $9=3^2$ δεν διαιρείται, γιατί το $3$ εμφανίζεται μόνο μία φορά.
\sol{12} $41,43,47,53,59$.
\sol{13} $(3,5),(5,7),(11,13),(17,19),(29,31),(41,43)$.
\sol{14} α) $8\cdot 3\cdot 7=168$. β) Ο $2\cdot 3\cdot 5\cdot 7=210$.
\sol{15} Π.χ. $33=3\cdot 11$, $63=7\cdot 9$, $93=3\cdot 31$ (και $143=11\cdot 13$).
\sol{16} Αν ένας σύνθετος αριθμός μέχρι το $100$ είχε όλους τους πρώτους παράγοντές του μεγαλύτερους ή ίσους του $11$, θα ήταν τουλάχιστον $11\cdot 11=121>100$. Άρα έχει πρώτο παράγοντα μικρότερο του $11$, δηλαδή έναν από τους $2,3,5,7$. Έτσι τα πολλαπλάσια των $2,3,5,7$ αρκούν.
\sol{17} $6.030$: με $2,3,5,9,10$ (άθροισμα $9$, όχι με το $4$ γιατί $30:4$ δεν είναι τέλεια). $2.800$: με $2,4,5,10$ (άθροισμα $10$, όχι με $3$ ή $9$). $4.617$: με $3,9$ (άθροισμα $18$), όχι με $2,4,5,10$.
\sol{18} $325=5\cdot 65=5\cdot 5\cdot 13=5^2\cdot 13$. $338=2\cdot 169=2\cdot 13^2$.
